\chapter{第二单元}

\Needspace{12\baselineskip}

\section{第二单元素养评价（A）}

\Needspace{8\baselineskip}

\pdfbookmark[2]{原书第9页}{orig-9}

\subsection*{原书第9页}

% p009_s01
\BookQuestion{一、1（逐组）}
下面的各组线中，（\hspace{1.8em}）互相平行；（\hspace{1.8em}）相交，（\hspace{1.8em}）互相垂直。（填序号）\par
\BookFigure{figures/p009_s01-01.tex}
\begin{BookAnswer}
①②⑦是两条不相交的直线，互相平行；④⑤⑥的直线延长后相交；⑥相交成直角，互相垂直。③画的是曲线，不按直线的平行或相交分类。因此三空依次填①②⑦、④⑤⑥、⑥。
\end{BookAnswer}

% p009_s02
\BookQuestion{一、2（逐图）}
下面哪个图形中，既有互相平行的线段，又有互相垂直的线段？在它下面画“\(\checkmark\)”。\par
\BookFigure{figures/p009_s02-01.tex}
\begin{BookAnswer}
只给第②个图形下面打\(\checkmark\)。①有平行边但无直角；②上、下底平行，左边又与底垂直，两种关系都有；③有平行边但无直角；④有直角但无平行边；⑤有互相平行的横线，但斜边不与横线垂直。
\BookFigure{figures/p009_s02-02.tex}
\end{BookAnswer}

% p009_s03
\BookQuestion{一、3（全部空）}
同一平面内，两条直线\(m\)和\(n\)的位置关系只有（\hspace{1.8em}）和（\hspace{1.8em}），如果直线\(m\)和\(n\)相交成直角，那么直线\(m\)和\(n\)互相（\hspace{1.8em}），记作（\hspace{1.8em}），读作（\hspace{1.8em}）。\par
\begin{BookAnswer}
同一平面内两条不同直线只有“平行”和“相交”两种位置关系。相交又分为“垂直”和“不垂直”，所以垂直属于相交。相交成直角时叫“互相垂直”，记作\(m\perp n\)，读作“\(m\)垂直于\(n\)”。
\end{BookAnswer}

% p009_s04
\BookQuestion{一、4}
从客厅的吊灯点到地面画三条线段，长度分别是3米、4米、5米，其中一条是垂线段，那么吊灯点到地面的距离是（\hspace{1.8em}）米。\par
\begin{BookAnswer}
点到直线的所有线段中，垂线段最短。\(3<4<5\)，所以垂线段长\(3\)米，距离填\(3\)。
\end{BookAnswer}

% p009_s05
\BookQuestion{一、5（全部空）}
长方形操场的长是100米，宽60米，操场的对边互相（\hspace{1.8em}），邻边互相（\hspace{1.8em}）。沿着跑道跑一圈，跑过的路程是（\hspace{1.8em}）米。\par
\begin{BookAnswer}
长方形对边平行、邻边垂直；沿跑道一圈就是周长：\((100+60)\times2=320\text{(米)}\)。依次填“平行、垂直、\(320\)”。
\end{BookAnswer}

% p009_s06
\BookQuestion{一、6（两个空）}
把一张长方形纸沿同一方向对折两次，展开后的折痕互相（\hspace{1.8em}），如果沿横竖不同方向对折两次，折痕互相（\hspace{1.8em}）。\par
\begin{BookAnswer}
沿同一方向对折两次，各条折痕互相\textbf{平行}；分别沿横向、竖向对折，两条折痕互相\textbf{垂直}。
\end{BookAnswer}

% p009_s07
\BookQuestion{一、7（跨页）}
平行的健身步道两边之间，能画（\hspace{1.8em}）条垂线段，这些线段的长度都（\hspace{1.8em}），这个性质在修路时保证了路面宽度一致。（第二空续于原书第10页。）\par
\begin{BookAnswer}
平行线之间可以在不同位置画\textbf{无数}条垂线段，这些垂线段都表示同一对平行线的距离，所以长度\textbf{相等}。两个空分别填“无数、相等”。
\end{BookAnswer}

\Needspace{8\baselineskip}

\pdfbookmark[2]{原书第10页}{orig-10}

\subsection*{原书第10页}

% p010_s01
\BookQuestion{一、7（续）}
（接原书第9页第7题）平行的健身步道两边之间垂线段的长度都（\hspace{1.8em}），这个性质在修路时保证了路面宽度一致。\par
\begin{BookAnswer}
填“相等”。一对平行线之间的垂直距离处处相等，路的两边平行就能保证宽度一致。
\end{BookAnswer}

% p010_s02
\BookQuestion{一、8（两个空）}
军军21时整准时睡觉，这时钟面上时针与分针所夹的角是（\hspace{1.8em}）度，我们也可以说时针与分针互相（\hspace{1.8em}）。\par
\begin{BookAnswer}
先换成钟面上的时刻：\[21\text{时}=\text{晚上}9\text{时}\]整点分针指向\(12\)，时针指向\(9\)，较小夹角是\(90^\circ\)，所以互相垂直。第一空后已有“度”，只填\(90\)；第二空填“垂直”。
\end{BookAnswer}

% p010_s03
\BookQuestion{二、1}
下列图形中，只有一组对边平行的是（\hspace{1.8em}）。\par
\noindent A.长方形\par
\noindent B.平行四边形\par
\noindent C.梯形\par
\noindent D.六边形\par
\begin{BookAnswer}
选\(\mathrm C\)。梯形只有一组对边平行；长方形、平行四边形都有两组对边平行。
\end{BookAnswer}

% p010_s04
\BookQuestion{二、2（读图条件）}
右图中\(C\)点到\(AB\)的距离是（\hspace{1.8em}）。\par
\noindent A.3厘米\par
\noindent B.4厘米\par
\noindent C.5厘米\par
\noindent D.不确定\par
\BookFigure{figures/p010_s04-01.tex}
\begin{BookAnswer}
选\(\mathrm A\)。线段\(AC\)与直线\(AB\)互相垂直，点\(C\)到直线\(AB\)的距离是垂线段\(AC\)的长度，即\(3\)厘米。
\end{BookAnswer}

% p010_s05
\BookQuestion{二、3}
下列说法正确的是（\hspace{1.8em}）。\par
\noindent A.垂直的两条直线不一定相交\par
\noindent B.平行线无限延长也不会相交\par
\noindent C.两条线段平行，它们的长度一定相等\par
\noindent D.以上说法都对\par
\begin{BookAnswer}
选\(\mathrm B\)。平行线无限延长仍不相交。垂直的直线一定相交，所以\(\mathrm A\)错；两线段平行并不要求等长，所以\(\mathrm C\)错。
\end{BookAnswer}

% p010_s06
\BookQuestion{二、4}
通过直线外一点可以画（\hspace{1.8em}）条直线与已知直线平行。\par
\noindent A.1\par
\noindent B.2\par
\noindent C.3\par
\noindent D.无数\par
\begin{BookAnswer}
选\(\mathrm A\)。在同一平面内，过直线外一点，有且只有\(1\)条直线与已知直线平行。
\end{BookAnswer}

% p010_s07
\BookQuestion{二、5}
直线行驶的汽车留下的车轮印迹（\hspace{1.8em}）。\par
\noindent A.互相平行\par
\noindent B.互相垂直\par
\noindent C.相交\par
\noindent D.互相平行或互相垂直\par
\begin{BookAnswer}
选\(\mathrm A\)。汽车沿直线前行时，两侧车轮的行进方向相同，轨迹不相交，互相平行。
\end{BookAnswer}

% p010_s08
\BookQuestion{二、6}
直线\(m\perp n\)，\allowbreak{}\(n\perp l\)，直线\(m\)和\(l\)的位置关系是（\hspace{1.8em}）。\par
\noindent A.互相平行\par
\noindent B.互相垂直\par
\noindent C.相交\par
\noindent D.不确定\par
\begin{BookAnswer}
选\(\mathrm A\)。同一平面内，两条不同直线都垂直于同一条直线时，这两条直线互相平行。题中\(m\perp n\)、\(n\perp l\)，也就是\(m\perp n\)、\(l\perp n\)，所以\(m\parallel l\)。
\end{BookAnswer}

% p010_s09
\BookQuestion{二、7}
下列物品中，两条边互相垂直的是（\hspace{1.8em}）。\par
\noindent A.红领巾\par
\noindent B.三角尺\par
\noindent C.平行四边形装饰画\par
\noindent D.六边形杯垫\par
\begin{BookAnswer}
选\(\mathrm B\)。三角尺有一个直角，组成这个直角的两条边互相垂直。其余图示物品的通常形状不含直角。
\end{BookAnswer}

% p010_s10
\BookQuestion{二、8}
下列关于过点\(A\)画已知直线\(l\)的垂线，步骤顺序正确的是（\hspace{1.8em}）。①把三角板一条直角边与直线\(l\)重合；②标出垂直符号；③沿着三角板另一条直角边画直线；④平移三角板，使直角顶点与点\(A\)重合。\par
\noindent A.①②③④\par
\noindent B.②①③④\par
\noindent C.①④③②\par
\noindent D.③①②④\par
\begin{BookAnswer}
选\(\mathrm C\)，顺序①④③②：先把一条直角边与\(l\)重合，保持方向平移至直角顶点到\(A\)，沿另一条直角边画线，最后标直角符号。
\end{BookAnswer}

% p010_s11
\BookQuestion{二、9}
在两条平行的步道之间修最短的人行道，该怎么画图纸？（\hspace{1.8em}）\par
\noindent A.画一条斜线段\par
\noindent B.画一条垂线段\par
\noindent C.画一条任意线段\par
\noindent D.不确定\par
\begin{BookAnswer}
选\(\mathrm B\)。两条平行步道之间，连接两边的垂线段最短，故画一条垂线段。
\end{BookAnswer}

\Needspace{8\baselineskip}

\pdfbookmark[2]{原书第11页}{orig-11}

\subsection*{原书第11页}

% p011_s01
\BookQuestion{二、10}
工程队在修建高速路时，有时需要将弯曲的道路改直来缩短路程，这样做的道理是（\hspace{1.8em}）。\par
\noindent A.两点确定一条直线\par
\noindent B.两点之间线段最短\par
\noindent C.过一点有无数条直线\par
\noindent D.不确定\par
\begin{BookAnswer}
选\(\mathrm B\)。相同的两个端点之间，线段比弯折的路径短；数学依据是“两点之间线段最短”。
\end{BookAnswer}

% p011_s02
\BookQuestion{三、1（1）}
列竖式计算：\(114\times23=\textcolor{answerred}{2622}\)。\par
\begin{BookAnswer}
把\(23\)拆成\(3+20\)，两个部分积分别为\(342\)和\(2280\)，合起来是\(2622\)。竖式各列按数位对齐：\[\begin{array}{r}114\\\times23\\\hline342\\2280\\\hline2622\end{array}\]
\end{BookAnswer}

% p011_s03
\BookQuestion{三、1（2）}
列竖式计算：\(448\div48=\textcolor{answerred}{9\cdots16}\)。\par
\begin{BookAnswer}
\[\begin{array}{r@{\;}r}&9\\\cline{2-2}48\big)&448\\&432\\\cline{2-2}&16\end{array}\]\(448\div48\)商\(9\)、余\(16\)。因为\(48\times9=432\)，\allowbreak{}\(448-432=16\)，且\(16<48\)。检验\(48\times9+16=448\)。
\end{BookAnswer}

% p011_s04
\BookQuestion{三、1（3）}
列竖式计算：\(372\div62=\textcolor{answerred}{6}\)。\par
\begin{BookAnswer}
\[\begin{array}{r@{\;}r}&6\\\cline{2-2}62\big)&372\\&372\\\cline{2-2}&0\end{array}\]\(372\div62=6\)。因为\(62\times6=372\)，正好除尽，余数为\(0\)。
\end{BookAnswer}

% p011_s05
\BookQuestion{三、1（4）}
列竖式计算：\(598\div18=\textcolor{answerred}{33\cdots4}\)。\par
\begin{BookAnswer}
\[\begin{array}{r@{\;}r}&33\\\cline{2-2}18\big)&598\\&54\phantom{0}\\\cline{2-2}&58\\&54\\\cline{2-2}&4\end{array}\]\(598\div18\)商\(33\)、余\(4\)。先用前两位\(59\)商\(3\)，余\(5\)；落下\(8\)成\(58\)，再商\(3\)，余\(4\)。检验\(18\times33+4=598\)，且\(4<18\)。
\end{BookAnswer}

% p011_s06
\BookQuestion{三、2（1）}
计算下列各题：\(1008\div[(73+32)\div5]\)。\par
\begin{BookAnswer}
\[\begin{aligned}&1008\div[(73+32)\div5]\\&=1008\div[105\div5]\\&=1008\div21\\&=48\end{aligned}\]
\end{BookAnswer}

% p011_s07
\BookQuestion{三、2（2）}
计算下列各题：\(966\div(27+19)\times45\)。\par
\begin{BookAnswer}
\[\begin{aligned}&966\div(27+19)\times45\\&=966\div46\times45\\&=21\times45\\&=945\end{aligned}\]\par\textbf{简便提示（学完运算律后）}\par \[\begin{aligned}&21\times45\\={}&(20+1)\times45\\={}&900+45\\={}&945\end{aligned}\]
\end{BookAnswer}

% p011_s08
\BookQuestion{三、2（3）}
计算下列各题：\(28\div4+63\times25\)。\par
\begin{BookAnswer}
\[\begin{aligned}&28\div4+63\times25\\&=7+1575\\&=1582\end{aligned}\]\par\textbf{运算律口算提示（学完运算律后）}\par 补充乘法部分的口算过程：把63分成60与3，用乘法分配律分别求积。前面的28除以4仍得7，再与1575相加为1582。\[\begin{aligned}&63\times25\\={}&(60+3)\times25\\={}&60\times25+3\times25\\={}&1500+75\\={}&1575\end{aligned}\]
\end{BookAnswer}

% p011_s09
\BookQuestion{四、1}
过\(A\)点画出已知直线\(l\)的平行线和垂线。\par
\BookFigure{figures/p011_s09-01.tex}
\begin{BookAnswer}
两条所画直线都必须经过黑点\(A\)。\par 画垂线：把三角板的一条直角边靠齐\(l\)，沿\(l\)移动，使另一条直角边经过\(A\)，再沿这条边画线。用直角检查交角，并把直角符号标在与\(l\)的交点处。\par 画平行线：三角板一边靠齐\(l\)，用直尺贴住另一边；固定直尺，沿尺平移三角板，移到作图边经过\(A\)时画线。移动时保持方向不变。
\BookFigure{figures/p011_s09-02.tex}
\end{BookAnswer}

% p011_s10
\BookQuestion{四、2}
补充两条边，使下面的图形变成一个长方形。\par
\BookFigure{figures/p011_s10-01.tex}
\begin{BookAnswer}
已有两条边相邻且垂直。分别过两条已有边的另一端点，作与对面已有边平行的线段，使它们相交闭合，就补成了长方形。核对：对边分别平行，四个角均为直角。
\BookFigure{figures/p011_s10-02.tex}
\end{BookAnswer}

\Needspace{8\baselineskip}

\pdfbookmark[2]{原书第12页}{orig-12}

\subsection*{原书第12页}

% p012_s01
\BookQuestion{四、3}
一只小蚂蚁在图中的点\(O\)处，它想快点爬到平行四边形的边线上，请你帮它设计一条最近的爬行路线，在图中画出来。\par
\BookFigure{figures/p012_s01-01.tex}
\begin{BookAnswer}
点到某一条边的最近路线是向这条边作垂线段；要到“四条边中任意一条”，还要比较这四条垂线段的长度。按图示，从\(O\)向上边\(AD\)作的垂线段最短。若把\(O\)视作上、下两平行边的正中位置，则向\(AD\)或\(BC\)作出的等长垂线段都可以。
\BookFigure{figures/p012_s01-02.tex}
\end{BookAnswer}

% p012_s02
\BookQuestion{五、1（1）}
一个长方形花园的周长是84米，其中一组对边之间的距离是12米。（1）另一组对边之间的距离是多少米？\par
\begin{BookAnswer}
长方形一组对边间的距离就是其中一条边长，已知一边长\(12\)米。\par \par\noindent\begin{minipage}{\linewidth}\textbf{方法一：用半周长求另一边}\par \[\begin{aligned}&84\div2-12\\&=42-12\\&=30\text{(米)}\end{aligned}\]答：另一组对边之间的距离是\(30\)米。\end{minipage}\par \par \par\noindent\begin{minipage}{\linewidth}\textbf{方法二：从周长中先扣两条已知边}\par \[\begin{aligned}&(84-12\times2)\div2\\&=(84-24)\div2\\&=60\div2\\&=30\text{(米)}\end{aligned}\]答：另一组对边之间的距离是\(30\)米。\end{minipage}\par
\end{BookAnswer}

% p012_s03
\BookQuestion{五、1（2）}
一个长方形花园的周长是84米，其中一组对边之间的距离是12米。（2）这个长方形花园的面积是多少平方米？\par
\begin{BookAnswer}
一组对边间距是宽\(12\)米，上一问已求出长\(30\)米。用“长方形面积\(=\)长\(\times\)宽”：\par \noindent\begin{minipage}{\linewidth}\textbf{用长乘宽求面积}\par \[\begin{aligned}&(84\div2-12)\times12\\&=(42-12)\times12\\&=30\times12\\&=360\text{(平方米)}\end{aligned}\]答：花园面积是\(360\)平方米。\end{minipage}\par 分步写成\(30\times12\)与上面的综合算式表达的是同一思路，不另算第二种解法。
\end{BookAnswer}

% p012_s04
\BookQuestion{五、2（1）}
2013年至今，国家出台并落实惠民政策，由各级政府出资对城镇道路进行升级改造，惠及千家万户。（1）计划分别从\(A\)、\(B\)两村出发各修一条公路与国道连接，怎样修能使修的路最短？请在图中画出来。\par
\BookFigure{figures/p012_s04-01.tex}
\begin{BookAnswer}
分别从\(A\)村和\(B\)村向国道作垂线段，连接村庄黑点与对应垂足。每条线段都应与国道边线垂直，直角符号标在垂足处。依据：点到直线的垂线段最短。
\BookFigure{figures/p012_s04-02.tex}
\end{BookAnswer}

% p012_s05
\BookQuestion{五、2（2）}
（接上题）（2）计划在\(A\)、\(B\)两村之间修一座高架桥，怎样修最短？请在图中画出来。\par
\BookFigure{figures/p012_s05-01.tex}
\begin{BookAnswer}
直接连接\(A\)、\(B\)两个村庄的黑点，所成线段就是最短的高架桥路线。这里比较的是两个固定端点之间的路线，依据“两点之间线段最短”。
\BookFigure{figures/p012_s05-02.tex}
\end{BookAnswer}

% p012_s06
\BookQuestion{五、3}
下面是某城市的局部交通线路图，计划再修一条阳光大道，规划为民主路的左侧经过点\(A\)与新华路平行，在民主路的右侧与青年路平行。请你画出阳光大道规划线路。\par
\BookFigure{figures/p012_s06-01.tex}
\begin{BookAnswer}
左段必须经过黑点\(A\)，并与新华路平行；把这段线延伸到民主路，转折点落在民主路上。再从同一个转折点画右段，使它与青年路平行。两段首尾共用一个端点，才是一条连续规划线。
\BookFigure{figures/p012_s06-02.tex}
\end{BookAnswer}

\Needspace{12\baselineskip}

\section{第二单元素养评价（B）}

\Needspace{8\baselineskip}

\pdfbookmark[2]{原书第13页}{orig-13}

\subsection*{原书第13页}

% p013_s01
\BookQuestion{一、1（两个空）}
如下图，互相平行的是（\hspace{1.8em}），互相垂直的是（\hspace{1.8em}）。\par
\BookFigure{figures/p013_s01-01.tex}
\begin{BookAnswer}
①③的两条直线方向相同，延长也不相交，互相平行；④两条直线相交成直角，互相垂直。故依次填①③、④。②相交但不垂直，⑤也不成直角，⑥是曲线。
\end{BookAnswer}

% p013_s02
\BookQuestion{一、2（全部空）}
直线\(m\)与直线\(n\)相交成直角时，这两条直线的位置关系是（\hspace{1.8em}），记作（\hspace{1.8em}），读作（\hspace{1.8em}）；这两条直线的交点叫作（\hspace{1.8em}）。\par
\begin{BookAnswer}
直线相交成直角，叫“互相垂直”，记作\(m\perp n\)，读作“\(m\)垂直于\(n\)”。这两条直线的交点叫“垂足”。
\end{BookAnswer}

% p013_s03
\BookQuestion{一、3}
钟表的圆形表盘上，除了3时整，（\hspace{1.8em}）时整的时针和分针所在直线也互相垂直。\par
\begin{BookAnswer}
填\(9\)。\(3\)时与\(9\)时的分针都指向\(12\)，时针分别指\(3\)、\(9\)，两种情况的较小夹角都是\(90^\circ\)。
\end{BookAnswer}

% p013_s04
\BookQuestion{一、4}
过春节贴春联是我国的传统习俗。如图，对联横批贴歪了，爷爷说：只要让横批的下边与门框的上边（\hspace{1.8em}），横批就正了。\par
\BookFigure{figures/p013_s04-01.tex}
\begin{BookAnswer}
填“平行”或“互相平行”。要把横批贴正，横批下边与门框上边的方向应一致，延长也不相交。
\end{BookAnswer}

% p013_s05
\BookQuestion{一、5（两个空）}
教室的推拉玻璃窗，推拉轨道的上下两条边互相（\hspace{1.8em}），玻璃窗的竖直边框和轨道边的位置关系是（\hspace{1.8em}）。\par
\begin{BookAnswer}
推拉轨道的上、下两边都是横向，互相\textbf{平行}；玻璃窗的竖直边框与横向轨道相交成直角，互相\textbf{垂直}。
\end{BookAnswer}

% p013_s06
\BookQuestion{一、6（全部空）}
如图，直线\(EG\)平行于直线\(FH\)，如果线段\(EF\)的长度是5厘米，那么线段\(GH\)的长度是（\hspace{1.8em}），线段\(AD\)的长度是（\hspace{1.8em}），它们的长度（\hspace{1.8em}），原因是（\hspace{1.8em}），线段\(AB\)与线段\(AC\)的长度一定（\hspace{1.8em}）5厘米（填“大于”“小于”或“等于”）。\par
\BookFigure{figures/p013_s06-01.tex}
\begin{BookAnswer}
由图中直角标记，\allowbreak{}\(EF,GH,AD\)都是\(EG\)与\(FH\)间的垂线段。一对平行线之间的距离处处相等，故\(GH=AD=EF=5\)厘米。前两个数值空后没有预印单位，应分别填\(5\text{厘米}\)、\(5\text{厘米}\)；第三空填“相等”；原因填“平行线间距离处处相等”。\par \(AD\)是\(A\)到\(FH\)的垂线段，\allowbreak{}\(AB,AC\)都是斜线段，故\(AB>AD,AC>AD\)，最后填“大于”。
\end{BookAnswer}

% p013_s07
\BookQuestion{一、7（两个空）}
在同一平面内，有4条直线，其中\(a\parallel b\)，\allowbreak{}\(b\perp c\)，\allowbreak{}\(a\)与\(c\)的位置关系是（\hspace{1.8em}），\allowbreak{}\(b\perp d\)，\allowbreak{}\(b\perp c\)，那么\(c\)与\(d\)的位置关系是（\hspace{1.8em}）。（第二空续于原书第14页。）\par
\BookFigure{figures/p013_s07-01.tex}
\begin{BookAnswer}
由\(a\parallel b\)且\(b\perp c\)，知\(a\perp c\)，第一空填“互相垂直”。\par 由\(b\perp d\)且\(b\perp c\)，同一平面内\(c,d\)都垂直于\(b\)，所以\(c\parallel d\)，第二空填“互相平行”（在下一页）。
\end{BookAnswer}

\Needspace{8\baselineskip}

\pdfbookmark[2]{原书第14页}{orig-14}

\subsection*{原书第14页}

% p014_s01
\BookQuestion{一、7（续）}
（接原书第13页第7题）\(b\perp d\)，\allowbreak{}\(b\perp c\)，那么\(c\)与\(d\)的位置关系是（\hspace{1.8em}）。\par
\begin{BookAnswer}
填“互相平行”。因为\(c\perp b,d\perp b\)，两条不同直线垂直于同一条直线时互相平行。
\end{BookAnswer}

% p014_s02
\BookQuestion{二、1}
在下面的字母中，线段的位置关系既有平行又有垂直的是（\hspace{1.8em}）。\par
\BookFigure{figures/p014_s02-01.tex}
\begin{BookAnswer}
选\(\mathrm B\)，字母\(\mathrm H\)的两条竖线平行，中间横线又与竖线垂直，兼有两种关系。
\end{BookAnswer}

% p014_s03
\BookQuestion{二、2}
我国战国时期《墨经》中记载“平，同高也”，它的意思是两条平行线间的距离（\hspace{1.8em}）。\par
\noindent A.不相等\par
\noindent B.处处相等\par
\noindent C.有时相等\par
\noindent D.无法确定\par
\begin{BookAnswer}
选\(\mathrm B\)。“平，同高也”说明平行线之间的垂直距离在各处相等，而不是有时相等。
\end{BookAnswer}

% p014_s04
\BookQuestion{二、3（两个空）}
在同一平面内，三条直线两两相交，最多有（\hspace{1.8em}）个交点，最少有（\hspace{1.8em}）个交点。\par
\noindent A.0\par
\noindent B.1\par
\noindent C.2\par
\noindent D.3\par
\begin{BookAnswer}
最多有\(3\)个交点，选\(\mathrm D\)：三条直线两两相交且不共点，形成三个不同交点。最少有\(1\)个交点，选\(\mathrm B\)：三条直线共同经过一点。题目已经要求“两两相交”，不能用平行线构造\(0\)个交点。
\end{BookAnswer}

% p014_s05
\BookQuestion{二、4}
若\(a\parallel b\)，\allowbreak{}\(b\parallel c\)，则\(a\)与\(c\)的位置关系是（\hspace{1.8em}）。\par
\noindent A.互相平行\par
\noindent B.相交\par
\noindent C.互相垂直\par
\noindent D.可能互相平行，也可能互相垂直\par
\begin{BookAnswer}
选\(\mathrm A\)。\(a\parallel b,b\parallel c\)，说明\(a,c\)都与同一直线\(b\)平行，因此\(a\parallel c\)。
\end{BookAnswer}

% p014_s06
\BookQuestion{二、5}
美术手工课上，同学们将一张正方形纸片对折2次，展开后得到的两条折痕（\hspace{1.8em}）。\par
\noindent A.互相垂直\par
\noindent B.可能互相垂直，也可能互相平行\par
\noindent C.互相平行\par
\noindent D.既不互相垂直，也不互相平行\par
\begin{BookAnswer}
按通常的折纸方式，选\(\mathrm B\)。两次沿同一方向对折，可得到平行折痕；分别沿横、竖方向对折，可得到垂直折痕。题目没有限定两次的方向，不能只选其中一种。
\end{BookAnswer}

% p014_s07
\BookQuestion{二、6}
某景区救援船停在人工湖的\(B\)点，发现\(A\)点有游客落水，救援船从\(B\)点赶往\(A\)点的最短行驶路线，对应的数学原理是（\hspace{1.8em}）。\par
\noindent A.垂直线段最短\par
\noindent B.两点之间线段最短\par
\noindent C.平行线间距离相等\par
\noindent D.过一点有且只有一条直线与已知直线平行\par
\BookFigure{figures/p014_s07-01.tex}
\begin{BookAnswer}
选\(\mathrm B\)。\(A,B\)是两个已确定的点，要走最短的救援路线，就沿线段\(AB\)直接前进；依据是“两点之间线段最短”。
\end{BookAnswer}

% p014_s08
\BookQuestion{二、7}
国家推进城乡供水一体化改造，在农村饮水安全工程中，从村庄垂直向河边拉输水管道，水管最短，节约管材、降低成本，其中蕴含的数学道理是（\hspace{1.8em}）。\par
\noindent A.两点之间线段最短\par
\noindent B.垂直线段最短\par
\noindent C.两点确定一条直线\par
\noindent D.平行线间距离相等\par
\begin{BookAnswer}
选\(\mathrm B\)。河边在题中看作一条直线，村庄是线外一点。从村庄垂直到河边的线段最短；依据是“垂线段最短”。
\end{BookAnswer}

\Needspace{8\baselineskip}

\pdfbookmark[2]{原书第15页}{orig-15}

\subsection*{原书第15页}

% p015_s01
\BookQuestion{三、1（1）}
列竖式计算：\(258\times13=\textcolor{answerred}{3354}\)。\par
\begin{BookAnswer}
\(258\times13\)的部分积是\(774\)与\(2580\)，合计\(3354\)。\[\begin{array}{r}258\\\times13\\\hline774\\2580\\\hline3354\end{array}\]
\end{BookAnswer}

% p015_s02
\BookQuestion{三、1（2）}
列竖式计算：\(754\div26=\textcolor{answerred}{29}\)。\par
\begin{BookAnswer}
\[\begin{array}{r@{\;}r}&29\\\cline{2-2}26\big)&754\\&52\phantom{0}\\\cline{2-2}&234\\&234\\\cline{2-2}&0\end{array}\]\(754\div26=29\)。先用\(75\)商\(2\)，\allowbreak{}\(26\times2=52\)，余\(23\)；落下\(4\)得\(234\)，再商\(9\)，\allowbreak{}\(26\times9=234\)，余\(0\)。检验\(26\times29=754\)。
\end{BookAnswer}

% p015_s03
\BookQuestion{三、1（3）}
列竖式计算：\(308\times25=\textcolor{answerred}{7700}\)。\par
\begin{BookAnswer}
\[\begin{array}{r}308\\\times25\\\hline1540\\6160\\\hline7700\end{array}\]\(308\times25\)的部分积是\(1540\)与\(6160\)，合计\(7700\)。也可把\(25\)看作\(100\div4\)，算\(308\div4\times100=77\times100=7700\)。
\end{BookAnswer}

% p015_s04
\BookQuestion{三、1（4）}
列竖式计算：\(552\div69=\textcolor{answerred}{8}\)。\par
\begin{BookAnswer}
\[\begin{array}{r@{\;}r}&8\\\cline{2-2}69\big)&552\\&552\\\cline{2-2}&0\end{array}\]\(552\div69=8\)。因为\(69\times8=552\)，正好除尽，余数是\(0\)。
\end{BookAnswer}

% p015_s05
\BookQuestion{三、2（1）}
计算下列各题：\(4000-15\times(24+195)\)。\par
\begin{BookAnswer}
\[\begin{aligned}&4000-15\times(24+195)\\&=4000-15\times219\\&=4000-3285\\&=715\end{aligned}\]\par\textbf{简便提示（学完运算律后）}\par \[\begin{aligned}&4000-15\times219\\={}&4000-15\times(220-1)\\={}&4000-(3300-15)\\={}&4000-3300+15\\={}&700+15\\={}&715\end{aligned}\]
\end{BookAnswer}

% p015_s06
\BookQuestion{三、2（2）}
计算下列各题：\(180\div[(84-12)\div4]\)。\par
\begin{BookAnswer}
\[\begin{aligned}&180\div[(84-12)\div4]\\&=180\div[72\div4]\\&=180\div18\\&=10\end{aligned}\]
\end{BookAnswer}

% p015_s07
\BookQuestion{三、2（3）}
计算下列各题：\(34\times[160-(60-15)]\)。\par
\begin{BookAnswer}
\[\begin{aligned}&34\times[160-(60-15)]\\&=34\times[160-45]\\&=34\times115\\&=3910\end{aligned}\]\par\textbf{简便提示（学完运算律后）}\par \[\begin{aligned}&34\times[160-(60-15)]\\={}&34\times[160-60+15]\\={}&34\times[100+15]\\={}&34\times100+34\times15\\={}&3400+510\\={}&3910\end{aligned}\]
\end{BookAnswer}

% p015_s08
\BookQuestion{三、2（4）}
计算下列各题：\([288\div(295-289)]\times17\)。\par
\begin{BookAnswer}
\[\begin{aligned}&[288\div(295-289)]\times17\\&=[288\div6]\times17\\&=48\times17\\&=816\end{aligned}\]\par\textbf{简便提示（学完运算律后）}\par 括号与除法照常求得48，最后可把48看成50少2，按乘法分配律同时算50乘17与2乘17。\[\begin{aligned}&48\times17\\={}&(50-2)\times17\\={}&50\times17-2\times17\\={}&850-34\\={}&816\end{aligned}\]
\end{BookAnswer}

% p015_s09
\BookQuestion{四、画一画}
过点\(A\)作已知直线的垂线段，在该垂线段上取点\(B\)，使\(AB=3\)厘米，再过点\(B\)画出已知直线的平行线。\par
\BookFigure{figures/p015_s09-01.tex}
\begin{BookAnswer}
先过黑点\(A\)作已知直线的垂线，再沿该垂线量取\(AB=3\)厘米，标出点\(B\)。最后过\(B\)作与原直线平行的直线。\par 核对三点：\(A,B\)都在同一条垂线上；\allowbreak{}\(AB\)的实际长度是\(3\)厘米；过\(B\)的新直线与原线方向相同。直角符号标在\(A\)处。
\BookFigure{figures/p015_s09-02.tex}
\end{BookAnswer}

% p015_s10
\BookQuestion{五、1}
小鸭很擅长游泳，在陆地上走路却很费劲。它打算从自己家到小鸡家玩耍（如图）。想一想，怎样走才能最省力气？请画出路线图。\par
\BookFigure{figures/p015_s10-01.tex}
\begin{BookAnswer}
题目强调小鸭在陆地上走路费劲，因此让两段\textbf{陆地路程最短}：从小鸭家向近岸作垂线段，到垂足下水；再从小鸡家向另一岸作垂线段，把垂足作为上岸点；水中连接两个垂足。\par 两段陆地路线都应与对应河岸垂直。这里依据“点到直线的垂线段最短”；水中连接两点则用线段最短。
\BookFigure{figures/p015_s10-02.tex}
\end{BookAnswer}

\Needspace{8\baselineskip}

\pdfbookmark[2]{原书第16页}{orig-16}

\subsection*{原书第16页}

% p016_s01
\BookQuestion{五、2（1）}
为了人身安全，过马路时要走人行横道。（1）如果丁丁从\(A\)点过马路，那么怎样走路线最短？请你在图中画一画。\par
\BookFigure{figures/p016_s01-01.tex}
\begin{BookAnswer}
从点\(A\)向马路对面作垂线段，沿人行横道直着走到对面。两侧道路边线平行，斜着过街会使路线变长；直角符号放在与对岸相交的垂足处。
\BookFigure{figures/p016_s01-02.tex}
\end{BookAnswer}

% p016_s02
\BookQuestion{五、2（2）}
（接上题）（2）如果丁丁的步行速度是每秒60厘米，按照最短路线行走，她过马路用了40秒，那么这条马路宽多少米？\par
\begin{BookAnswer}
先按“每秒\(60\)厘米、共\(40\)秒”求总长，再把厘米数换成米。\[100\text{厘米}=1\text{米}\]\par\noindent\begin{minipage}{\linewidth}\textbf{方法一：先求厘米数，再换算}\par \[\begin{aligned}&60\times40\\={}&2400\text{(厘米)}\end{aligned}\]\[2400\text{厘米}=(2400\div100)\text{米}=24\text{米}\]答：这条马路宽\(24\)米。\end{minipage}\par\noindent\begin{minipage}{\linewidth}\textbf{方法二：每十秒分为一段}\par 先同时求每段的厘米数和段数：\[60\times10=600\text{(厘米)},\qquad40\div10=4\text{(段)}\]\[600\text{厘米}=(600\div100)\text{米}=6\text{米}\]再用每段长度乘段数：\[\begin{aligned}&6\times4\\={}&24\text{(米)}\end{aligned}\]答：这条马路宽\(24\)米。\end{minipage}\par 每\(10\)秒走\(6\)米，\allowbreak{}\(40\)秒有这样的\(4\)段。\par 用相同的速度单位检验，先把米换成厘米：\[24\text{米}=2400\text{厘米}\]\[\begin{aligned}&2400\div40\\={}&60\text{(厘米/秒)}\end{aligned}\]还原题目每秒\(60\)厘米的速度。
\end{BookAnswer}

% p016_s03
\BookQuestion{五、3}
在下面的平行线之间画一个最大的正方形。\par
\BookFigure{figures/p016_s03-01.tex}
\begin{BookAnswer}
先作出两条平行线之间的一条垂线段，把它的长度作为正方形的边长。再取同样长的上、下两条边，并用另一条垂线段连接相应端点。\par 核对：上、下边分别落在给定平行线上；两条侧边都与平行线垂直；四边等长。边长等于平行线间距时达到最大。
\BookFigure{figures/p016_s03-02.tex}
\end{BookAnswer}

% p016_s04
\BookQuestion{六、动脑筋}
蓝蓝同学参加了一个户外探险活动，根据线索先找到藏宝地者获胜。请你根据以下线索找到藏宝地，并在图中画一画。线索一：有一条笔直的公路经过\(A\)点和\(B\)点。线索二：藏宝地与\(C\)点的连线与公路平行。线索三：藏宝地在\(C\)点的左下方，与\(C\)点的距离等于\(C\)点到公路的距离。\par
\BookFigure{figures/p016_s04-01.tex}
\begin{BookAnswer}
①连接黑点\(A,B\)并向两端延长，这条直线表示笔直公路。\par ②从\(C\)向公路作垂线段，记垂足为\(D\)，线段\(CD\)表示\(C\)到公路的距离。\par ③过\(C\)作公路的平行线，在\(C\)的左下方截取\(CP=CD\)。点\(P\)就是藏宝地。\par 验证三条线索：公路经过\(A,B\)；\allowbreak{}\(CP\parallel AB\)；\allowbreak{}\(P\)在\(C\)左下方，且\(CP\)等于垂直距离\(CD\)。
\BookFigure{figures/p016_s04-02.tex}
\end{BookAnswer}
